\section{DiT-Air：参数共享能走多远}

在 MMDiT 走向更复杂之后，DiT-Air 反向提问：哪些参数真的需要每层、每模态独立？共享这些调制、投影或前馈模块能降低成本，但共享过多会削弱模态专属能力。这个问题比“是否使用 MMDiT”更接近工程 Pareto 优化。

\subsection{从设计问题到共享架构}

参数共享有三个维度：跨层共享、跨模态共享、跨子模块共享。不同维度对容量和运行成本影响不同。下面四页给出从问题、结构到 ablation 的完整链。本节阅读时应始终分开三个口径：权重是否减少、每次前向的 FLOPs 是否减少、真实硬件上的 latency 是否减少；只有三者同时核算，共享架构才具有工程意义。

\lecturefigure{slide-55.jpg}{简化问题：AdaLN、QKVO、MLP 与 attention 形式能共享多少}{55}

PixArt 已证明共享 AdaLN 有效；下一步是共享 Q/K/V/O projections 或 MLP，甚至用统一 joint attention 代替 self+cross stacks。每项共享减少参数，但不一定减少 activation 或 attention FLOPs。

\lecturefigure{slide-56.jpg}{DiT-Air 架构：以统一 block 比较不同参数共享位置}{56}

图中将 baseline、AdaLN sharing、QKVO sharing 与 MLP sharing 并列。统一表示便于做 controlled ablation。设计时要区分“参数少”“FLOPs 少”“壁钟快”：共享权重可能减少存储，却不减少每个 token 的矩阵乘法次数。

\lecturefigure{slide-57.jpg}{DiT-Air 的共享实验：不同组合形成质量--效率曲线}{57}

曲线显示部分共享可在较小质量损失下提升效率。读图时先找 Pareto frontier，再看误差条和计算口径。若共享导致 prompt fidelity 或细粒度文本下降，平均 FID 未必能显示问题。

\lecturefigure{slide-58.jpg}{Ablation 表：共享 QKVO/MLP 的收益与质量折损}{58}

表格把组件级选择变为证据。结论不是“共享越多越好”，而是 AdaLN sharing 往往较安全，QKVO/MLP sharing 更依赖目标与预算。面向低延迟部署可接受小幅质量损失，研究上追求上限则可能保留专属参数。

\teachervoice{01:05:00--01:08:43，老师用 DiT-Air 回答“是否真的需要 MMDiT”：不一定。AdaLN、QKVO 和 MLP 都可以共享，但效率收益必须与质量折损一起看。}

\begin{importantbox}{三种成本不要混为一谈}
Parameter count 影响权重存储；FLOPs 近似计算量；latency 还受 kernel、并行度和内存带宽影响。共享参数可能显著减小模型文件，却不一定按相同比例降低推理时间。
\end{importantbox}

\begin{lstlisting}[language=Python,caption={用 Pareto 指标比较共享方案，而不是只排单一质量分数}]
for design in candidate_architectures:
    quality = evaluate_prompt_fidelity_and_preference(design)
    memory = peak_device_memory(design)
    latency = benchmark_end_to_end_sampling(design)
    flops = estimate_denoiser_flops(design)
    report_pareto_point(design, quality, memory, latency, flops)
\end{lstlisting}

\subsection{本章小结}

DiT-Air 把架构简化转为可测的共享实验。AdaLN、QKVO 与 MLP sharing 分别影响参数、计算和容量；MMDiT 不是默认答案。工程选择应基于端到端 Pareto frontier，并用文本遵循、人类偏好和延迟共同评估。

